% Source: source/普通高考/2014/2014天津文.pdf, page 1, question 11.
% Flowchart topology: 开始 -> S=0,n=3 -> S=S+(-2)^n -> n=n-1 -> n<=1?;
% 否 loops to the entry above the first process, 是 -> 输出 S -> 结束.
\begin{tikzpicture}[
  x=1cm, y=0.86cm, >=Stealth,
  line width=0.45pt, line cap=round, line join=round,
  every node/.style={anchor=center, align=center,
    execute at begin node={\setlength{\parindent}{0pt}\noindent}},
  flowbox/.style={draw, minimum height=0.68cm, inner xsep=3pt, inner ysep=2pt,
    align=center, execute at begin node={\setlength{\parindent}{0pt}\noindent}},
  terminal/.style={flowbox, rounded corners=0.15cm, minimum width=1.20cm,
    minimum height=0.68cm},
  process/.style={flowbox, minimum width=2.85cm},
  decision/.style={flowbox, diamond, aspect=2.8, minimum width=3.25cm,
    minimum height=0.98cm, inner sep=1pt},
  io/.style={flowbox, trapezium, trapezium left angle=72, trapezium right angle=108,
    minimum width=1.90cm, minimum height=0.72cm},
  branchlabel/.style={anchor=center, align=center, inner sep=0pt,
    execute at begin node={\setlength{\parindent}{0pt}\noindent}}
]
  \node[terminal] (start) at (0,10.00) {开始};
  \node[process, minimum width=2.55cm] (init) at (0,8.65) {$S=0,n=3$};
  \node[process, minimum width=3.45cm] (update) at (0,7.15) {$S=S+(-2)^n$};
  \node[process, minimum width=2.40cm] (dec) at (0,5.70) {$n=n-1$};
  \node[decision] (test) at (0,4.20) {$n\leqslant 1?$};
  \node[io] (output) at (0,2.65) {输出 $S$};
  \node[terminal] (finish) at (0,1.25) {结束};

  \draw[->] (start.south) -- (init.north);
  \coordinate (loopJoin) at (0,8.00);
  \draw (init.south) -- (loopJoin);
  \draw[->] (loopJoin) -- (update.north);
  \draw[->] (update.south) -- (dec.north);
  \draw[->] (dec.south) -- (test.north);
  \draw[->] (test.south) -- node[branchlabel, right=2pt] {是} (output.north);
  \draw[->] (output.south) -- (finish.north);

  \coordinate (loopLeft) at (-2.75,4.20);
  \coordinate (loopTop) at (-2.75,8.00);
  \draw (test.west) -- node[branchlabel, above=2pt] {否} (loopLeft) -- (loopTop);
  \draw (loopTop) -- (loopJoin);
\end{tikzpicture}
